\chapter*{Conclusion générale et perspectives}
\addcontentsline{toc}{chapter}{Conclusion générale et perspectives}
\markboth{Conclusion générale}{}

Ce projet de fin d'études avait pour ambition de concevoir une plateforme
décisionnelle complète pour l'analyse financière de la \sfbt. L'objectif a été
atteint : à partir d'états financiers bruts et imparfaits, une chaîne de
traitement automatisée produit aujourd'hui un diagnostic financier fiable,
comparatif et prédictif.

\medskip
Sur le plan des \textbf{réalisations}, le travail a permis de : (i) unifier et
fiabiliser quatre sources hétérogènes en un jeu de données cohérent de plus de
dix mille lignes ; (ii) enrichir l'historique de la \sfbt par la génération des
clôtures trimestrielles manquantes ; (iii) implémenter \textbf{46 ratios
financiers} conformes à un référentiel normalisé et les synthétiser en un
\textbf{score global de performance} (83,5/100 pour la \sfbt, en tête du secteur) ;
(iv) structurer les données dans un \textbf{entrepôt en étoile} interrogeable ;
et (v) développer un module de \textbf{prévision} comparant deux modèles
d'apprentissage automatique, avec une précision de 85 à 96~\%. L'ensemble est
\textbf{testé} (70~tests) et \textbf{documenté}.

\medskip
Sur le plan des \textbf{apports}, ce projet illustre la mise en œuvre concrète de
la chaîne décisionnelle enseignée en \textit{Business Analytics}, depuis la
fiabilisation des données jusqu'à la modélisation prédictive, sur un cas réel
d'entreprise. Il met également en évidence l'importance de la qualité des données
et de la transparence sur leurs limites.

\medskip
Plusieurs \textbf{perspectives} d'amélioration se dégagent :
\begin{itemize}[noitemsep]
  \item \textbf{compléter les données manquantes} de la \sfbt (exercices annuels
        2021 et 2022) pour une couverture intégrale ;
  \item \textbf{finaliser les tableaux de bord Power BI} et, éventuellement, une
        interface web interactive ;
  \item \textbf{affiner la prévision} en intégrant la saisonnalité réelle et en
        testant des modèles de séries temporelles dédiés (SARIMA, Prophet) ;
  \item \textbf{enrichir le benchmark} avec d'autres acteurs du secteur et
        homogénéiser les devises pour comparer aussi les montants absolus ;
  \item \textbf{étendre la solution} à d'autres sociétés, la modularité de
        l'architecture le permettant aisément.
\end{itemize}

\medskip
En définitive, ce travail constitue une base solide, extensible et reproductible
pour le pilotage de la performance financière, et ouvre la voie à des analyses
décisionnelles plus poussées.
